\documentclass[10pt,a4paper]{article}
\usepackage[T1]{fontenc}
\usepackage[utf8]{inputenc}
\usepackage{lmodern}
\usepackage{amsmath,amssymb}
\usepackage[margin=21mm]{geometry}
\usepackage{microtype}
\usepackage[hidelinks]{hyperref}
\usepackage{enumitem}
\setlist[itemize]{leftmargin=*,nosep,topsep=3pt}
\setlength{\parindent}{0pt}
\setlength{\parskip}{5pt}
\setlength{\emergencystretch}{3em}
\newcommand{\question}[1]{\par\medskip{\normalsize\bfseries #1}\par\smallskip}
\pagestyle{plain}

\begin{document}
{\Large\bfseries Dopamine and neuromodulation accelerate the neural dynamics of volitional action in Parkinson's disease\par}

\textbf{Paper:} Koehler et al. (2023).\quad bioRxiv preprint; overlaps note 04.\par
\href{https://doi.org/10.1101/2023.10.30.564700}{DOI: \nolinkurl{10.1101/2023.10.30.564700}}\par
{\footnotesize Page references below use PDF page numbers, including front matter.}

\question{To which general question does this paper relate?}
Movement-intention measurements and a possible treatment objective based on successful action initiation. Provides the preprint version of the study in note~04.

\question{How does it answer the general questions?}
\textbf{Measured.} Combined cortical and STN recordings in 25 patients were analysed with readiness potentials, movement-intention classifiers and directional connectivity measures. EMG provided the execution reference (PDF pp.~1--8).

\textbf{Finding.} Medication and STN-DBS shortened intention-to-execution latencies, accompanied by a shift from beta towards theta cortico-subthalamic coupling. This motivates studying restoration of movement-related network dynamics. The preprint does not demonstrate that directly suppressing beta alone is the therapeutic mechanism (PDF pp.~8--12).

\textbf{Model relevance.} Supplies observable signals and behavioural timing measures, but no identified delayed neural rate model, actuator map, or stability certificate. The measured preparation interval is not an identified neural transmission delay.

\textbf{Version.} The local PDF is explicitly a non-peer-reviewed preprint. Use the published paper in note~04 for the final account and count the shared experiment once in a synthesis.

\question{What papers does it cite that could provide more information?}
\begin{itemize}
\item Lofredi et al. (2018), \emph{Dopamine-dependent scaling of subthalamic gamma bursts with movement velocity in patients with Parkinson's disease} (ref.~39, PDF p.~20).
\item Gilron et al. (2020), \emph{Chronic wireless streaming of invasive neural recordings at home for circuit discovery and adaptive stimulation}: cited in its preprint form (ref.~40, PDF p.~20).
\end{itemize}

\end{document}
